% Paper 2: Bitcoin as Unique Neutral Settlement: A Seven-Property Elimination
% Compiled from BGT-PAPER-2.md via pandoc + template assembly
% Build: xelatex BGT-PAPER-2.tex

\documentclass[12pt]{article}

% --- Fonts and encoding ---
\usepackage{fontspec}
\setmainfont{Latin Modern Roman}
\usepackage{amsmath}
% Note: amssymb omitted — unicode-math provides all symbols for XeLaTeX
\usepackage{unicode-math}
\setmathfont{Latin Modern Math}

% --- Page layout ---
\usepackage[margin=1in]{geometry}
\usepackage{setspace}
\onehalfspacing

% --- Theorem environments ---
\usepackage{amsthm}
\newtheorem{theorem}{Theorem}
\newtheorem{proposition}{Proposition}
\newtheorem{lemma}{Lemma}
\newtheorem{corollary}{Corollary}
\theoremstyle{definition}
\newtheorem{definition}{Definition}
\newtheorem{assumption}{Assumption}
\theoremstyle{remark}
\newtheorem*{remark}{Remark}
\renewcommand{\qedsymbol}{}

% --- Tables ---
\usepackage{booktabs}
\usepackage{longtable}
\usepackage{array}
\usepackage{tabularx}

% --- Page numbers ---
\usepackage{fancyhdr}
\pagestyle{fancy}
\fancyhf{}
\fancyfoot[C]{\thepage}
\renewcommand{\headrulewidth}{0pt}

% --- Misc ---
\usepackage{enumitem}
\usepackage{hyperref}
\hypersetup{
  colorlinks=true, linkcolor=black, citecolor=black, urlcolor=blue,
  pdftitle={Bitcoin as Unique Neutral Settlement: A Seven-Property Elimination},
  pdfauthor={Sean Hash},
  pdfkeywords={Bitcoin, neutral settlement, asset comparison, digital scarcity, reserve assets, capital allocation, property elimination, game theory, network effects}
}
\def\UrlBreaks{\do\/\do-\do.\do=\do_\do?\do\&\do\%\do\a\do\b\do\c\do\d\do\e\do\f\do\g\do\h\do\i\do\j\do\k\do\l\do\m\do\n\do\o\do\p\do\q\do\r\do\s\do\t\do\u\do\v\do\w\do\x\do\y\do\z\do\A\do\B\do\C\do\D\do\E\do\F\do\G\do\H\do\I\do\J\do\K\do\L\do\M\do\N\do\O\do\P\do\Q\do\R\do\S\do\T\do\U\do\V\do\W\do\X\do\Y\do\Z}
\usepackage{microtype}

% --- Unicode character fallbacks (Greek + math symbols in text mode) ---
\usepackage{newunicodechar}
\newunicodechar{δ}{$\delta$}
\newunicodechar{ε}{$\varepsilon$}
\newunicodechar{λ}{$\lambda$}
\newunicodechar{σ}{$\sigma$}
\newunicodechar{∈}{$\in$}
\newunicodechar{∎}{\rule{0.5em}{0.5em}}
\newunicodechar{≈}{$\approx$}
\newunicodechar{≠}{$\neq$}
\newunicodechar{≫}{$\gg$}

% --- Suppress section numbering (manual numbers in headings) ---
\setcounter{secnumdepth}{0}

% --- Custom commands for readability ---
\newcommand{\E}{\mathbb{E}}

\begin{document}

% ============================================================
% TITLE BLOCK
% ============================================================

\begin{center}
{\LARGE\bfseries Bitcoin as Unique Neutral Settlement: A Seven-Property Elimination}

\bigskip

{\large Sean Hash}\\
\url{bitcoingametheory.com}\\
\href{mailto:sean@bitcoingametheory.com}{sean@bitcoingametheory.com}

\bigskip

April 2026\\[4pt]
{\normalsize Working Paper. Cite as Hash (2026b).}

\medskip

{\small
\textbf{JEL Codes:} G11, G12, G15, E42, E51, L14\\[4pt]
\textbf{Keywords:} Bitcoin, neutral settlement, asset comparison, digital scarcity, reserve assets, capital allocation, property elimination, game theory, network effects
}
\end{center}

\bigskip

% ============================================================
% ABSTRACT
% ============================================================

\begin{abstract}
\noindent
Which asset can serve as neutral settlement in a multipolar world? Hash
(2026a) establishes that the payoff advantage of Exit over Stay is
monotone increasing in adoption, and derives seven necessary
properties: protocol security, neutrality, permissionless access,
cheap finality, absolute scarcity, informational security, and
adaptive resilience. This paper applies those properties. We conduct a
systematic elimination across seven asset classes: fiat currencies,
sovereign bonds, equities, real estate, gold, alternative blockchain
tokens, and Bitcoin. Six fail at least one required property. Bitcoin
satisfies all seven because protocol neutrality is independent of
ownership concentration. Under proof-of-work with cryptographic
custody, holding coins confers no consensus vote, no ability to alter
supply, no ability to censor transactions, and no ability to freeze
other holders' keys. Capture surfaces at ETF wrappers, exchanges,
mining pools, and corporate treasuries are real, but each is
unilaterally exitable without counterparty consent, and even a
coalition controlling a majority of transferable supply could not
prevent anyone holding keys from settling. No other asset separates
ownership from protocol authority in this way. The analysis does not
require fiat collapse. Bitcoin operates as settlement finality layer
in parallel with fiat payment rails.
\end{abstract}

\newpage

\section{1. Introduction}

The Exit Game (Hash, 2026a) proves that the payoff advantage of moving
capital from capturable settlement systems to neutral settlement is
monotonically increasing in adoption. But the framework is agnostic
about \emph{which} asset serves as the Exit destination. It derives
seven necessary properties (P1--P7) and shows that the Exit advantage is
monotone increasing. The identity of the neutral settlement asset is an
empirical question, not a theoretical one.

This paper answers that question. We take P1--P7 as given and ask: which
existing asset satisfies all seven simultaneously? The answer is reached
by elimination, not by advocacy. Each asset class is tested against each
property. Six classes fail. One survives.

A second contribution distinguishes this paper from Hash (2026a). The
Exit Game treats capture as a binary, protocol-layer constraint. This
paper relaxes that idealization. Every asset has capture surfaces at
some layer, whether issuance, custody, venue, or physical possession.
The relevant question is not whether capture exists but whether any
feasible coalition can use that capture to alter settlement rules or
prevent key-holders from transacting. We formalize this as
Ownership-Protocol Independence (Proposition 3) and show that Bitcoin
is the unique asset satisfying it. The argument is therefore
comparative. Bitcoin is not uncapturable at every layer; its
application-layer captures cannot translate into protocol-layer control
under proof-of-work, because coin ownership confers no consensus
authority.

The structure is deliberate. By separating the game-theoretic argument
in Hash (2026a) from the asset identification in this paper, we make the
logic auditable: a reader who rejects our claim about Bitcoin need only
show that some alternative satisfies P1--P7, without engaging the game
theory. A reader who accepts Bitcoin's properties need only verify the
elimination table, without reconstructing the dominance proof.
\section{2. The Property Framework}

We summarize the seven properties derived in Hash (2026a), Section 3.3.
Each property blocks a specific attack class that would violate neutral
settlement, defined by Hash (2026a) as immunity from seizure,
debasement, and political capture.

\subsection{2.1 Formal Specifications}

\noindent\textbf{P1 (Protocol Security).} Mining constitutes a Nash equilibrium;
attack cost exceeds expected gain. Formally:
$C_\text{attack} + P_\text{collapse} > \Delta_\text{attack}$, and
$\sum \delta^t \cdot \E[\Pi_\text{honest}] > \Delta_\text{attack} - C_\text{attack}$
(Budish, 2018; Biais et al., 2019).

\noindent\textbf{P2 (Neutrality).} No issuer, no foundation, no governance
mechanism capable of unilateral rule changes. $|CS_\text{protocol}| = 0$.

\noindent\textbf{P3 (Permissionless Access).} Any agent can initiate settlement
without third-party permission. $\Pr(\text{censor access} \mid \text{scale}) \approx 0$.

\noindent\textbf{P4 (Cheap Finality).} Settlement of \$1 billion for fees that
scale with transaction bytes, not value: under \$1 at prevailing mempool
rates (1 sat/vB, April 2026) and under \$520 even at the highest
sustained congestion on record (\textasciitilde1,200 sat/vB during the
April 2024 halving; \textasciitilde500 sat/vB during the May 2023
BRC-20 event), with mathematical finality in under 60 minutes.

\noindent\textbf{P5 (Absolute Scarcity).} Fixed supply with zero supply
elasticity: $dS/dP = 0$.

\noindent\textbf{P6 (Informational Security).} Custody is mathematical (private
key knowledge). Cost of seizure scales superlinearly with targets.

\noindent\textbf{P7 (Adaptive Resilience).} Protocol upgrades via consensus
without introducing governance capture.

\subsection{2.2 Necessity and Sufficiency}

\begin{proposition}[Necessity]
P1--P7 are individually necessary. Removing any single property enables
the corresponding attack vector.
\end{proposition}

\noindent Proved in Hash (2026a).

\begin{proposition}[Sufficiency]
P1--P7 are jointly sufficient. Candidate additional properties either
reduce to P1--P7 or introduce non-structural criteria that would create
capture surface, violating P2.
\end{proposition}

\noindent For example, ``privacy'' as a candidate property reduces to P6
(informational security).

\subsection{2.3 Ownership-Protocol Independence}

The seven properties in Hash (2026a) are stated as binary constraints
at the protocol layer. At the application layer, every asset admits
capture surfaces that P1--P7 do not directly address. We introduce a
formal complement.

\begin{definition}[Protocol Capture Fraction]
For asset $A$, let $\kappa_\text{protocol}(A) \in [0, 1]$ denote the
fraction of $A$'s settlement rules, supply schedule, and transaction
validity that any feasible coalition can unilaterally alter.
\end{definition}

\begin{definition}[Ownership Capture Fraction]
Let $\kappa_\text{owner}(A) \in [0, 1]$ denote the maximum fraction of
$A$'s transferable supply any single coalition of state actors,
custodians, issuers, or venue operators can hold, seize, freeze, or
render non-transferable.
\end{definition}

\begin{proposition}[Ownership-Protocol Independence]
Neutral settlement requires $\kappa_\text{protocol}(A) = 0$. For assets
where consensus authority or governance rights are conferred by
ownership, $\kappa_\text{owner} \to 1$ implies $\kappa_\text{protocol}
\to 1$: equities via shares, proof-of-stake tokens via stake,
governance tokens via votes, stablecoins via issuer control. For
proof-of-work Bitcoin with cryptographic custody, the implication
fails: $\kappa_\text{protocol}(\text{Bitcoin}) = 0$ for any value of
$\kappa_\text{owner}(\text{Bitcoin}) \in [0, 1]$.
\end{proposition}

\noindent The proposition is a necessary complement to P1--P7. An asset
satisfying all seven protocol-layer properties still fails as neutral
settlement if ownership concentration translates into consensus or
issuance control. Conversely, an asset with high ownership
concentration remains neutral if protocol rules are enforced
independently of coin ownership. Under proof-of-work, consensus is
determined by hashrate and node operation, neither of which is
conferred by holding coins. A coalition holding 51\% of transferable
supply cannot alter the 21-million cap, cannot censor transactions it
does not see, cannot seize coins held by other key-holders, and cannot
change consensus rules without independent agreement from miners and
nodes. The elimination in \S3 applies P1--P7 first, then verifies
$\kappa_\text{protocol}(A) = 0$ for the survivor.

\section{3. Elimination}

\subsection{3.1 Method}

For each asset class, we identify which properties are violated and the
mechanism of violation. An asset fails the test if it violates
\emph{any} single property, since P1--P7 are individually necessary.
Where ownership concentration translates directly into protocol
control, we note it explicitly, since Proposition 3 makes this the
critical failure mode for governance-token and equity-like assets.

\subsection{3.2 Results}

Table 1 summarizes the outcome of the elimination. Six of seven asset
classes fail at least one necessary property; Bitcoin is the only
candidate that violates none. The buckets can also be read as points
along a capturability continuum: fiat and sovereign bonds at the
most-capturable end (issuer-controlled, denomination-controlled);
equities and real estate next (state-regulated corporate form, or
immovable and title-registered); gold and alt-L1 tokens further out
(portable-but-physical, or digital-but-governance-captured); and
Bitcoin at the non-capturable extreme (informational custody, no
issuer, no founder). The detailed analysis in \S3.3 works through
each row in turn.

\begin{table}[h]
\centering
\begin{tabularx}{\textwidth}{llX}
\toprule
Asset Class & Properties Violated & Mechanism \\
\midrule
C1a: Fiat currencies & P2, P5 & Central-bank issuer with discretionary supply schedule; $|CS_\text{protocol}| > 0$ by construction \\
C1b: CBDCs & P2, P3, P5, P6 & Fiat failures plus programmable censorship and state surveillance \\
C2: Sovereign bonds & P5 & Negative real yield under financial repression; denominated in debaseable currency \\
C3: Equities & P2 & Regulatory capture via margin requirements, trading halts, and dilution; ownership above 50\% of shares confers board control \\
C4: Real estate & P3, P4, P6 & Registration-gated transfers, multi-week settlement, physically seizable, subject to property tax and eminent domain \\
C5: Gold & P4, P6 & Settlement friction of 3--8\% for insurance, transport, and assay; physically seizable \\
C6a: Alt-L1 tokens & P2 & Governance capture via foundation treasuries, VC voting blocs, or proof-of-stake concentration \\
C6b: Stablecoins & P2, P5 & Issuer freeze function; denominated in fiat (inherits P5 debasement) \\
C7: Bitcoin & None identified & None \\
\bottomrule
\end{tabularx}
\end{table}

\subsection{3.3 Detailed Analysis}

\noindent\textbf{C1: Fiat Currencies and CBDCs.} Central banks set
monetary policy to serve domestic objectives (Hash, 2026a, Assumption
1). The supply schedule is discretionary: $dS/dP \neq 0$ by design.
This violates P5 directly. The central bank is itself the issuer with
unilateral authority to alter the supply schedule, which makes
$|CS_\text{protocol,\,fiat}| > 0$ and violates P2. CBDCs inherit both
failures and add programmable censorship (violating P3) and state
surveillance (violating P6). No fiat currency has maintained purchasing
power over a 50-year horizon against hard assets.

\noindent\textbf{C2: Sovereign Bonds.} Bonds are claims on future fiat. When the
underlying currency is debased, bondholders bear the loss through
negative real yields. Financial repression, holding interest rates below
inflation, is the observed mechanism. This is not a market failure; it
is policy working as intended. Real yields on 10-year US Treasuries
were negative for 26 of 36 months from 2020--2023 (Federal Reserve Bank
of St.~Louis, 2024, FRED series DGS10 minus T10YIE).

\noindent\textbf{C3: Equities.} Productive assets generate real returns but are
subject to regulatory capture. Governments can impose capital controls,
windfall taxes, forced sales, trading halts, and beneficial ownership
reporting requirements. The capture surface is large:
$|CS_\text{protocol,\,equities}| \gg 0$, and ownership concentration
above 50\% of shares confers board control, collapsing the ownership
and protocol layers into a single coalition.

\noindent\textbf{C4: Real Estate.} Property is the paradigmatic \emph{seizable}
asset. Eminent domain, property tax, and zoning regulation constitute
permanent capture surfaces. Physical location is known (violating P6).
Settlement requires title search, escrow, and legal process (violating
P4). Cross-border real estate settlement is measured in weeks, not
minutes.

\noindent\textbf{C5: Gold.} This is the hardest case and deserves extended
treatment. Gold stands here for the broader class of energy-to-produce
commodities (silver, platinum-group metals, industrial metals, and
historically also shells, salt, and cattle); gold is the representative
because it is the most durable and has the lowest and most predictable
supply growth rate within the class. Real estate (C4) and art are also commodity-like: each requires energy
or labor to produce and holds physical form. They differ from gold on
portability. Real estate is immovable, which makes it maximally
capturable by the political system that controls its jurisdiction;
this is why it separates out as its own bucket with a harder P6 failure
and an additional P3 failure via registration-gated transfers. Art is
portable but heterogeneous and provenance-dependent, which raises
verification cost above gold's and pushes P4 failures further.
Everything argued below about gold applies \emph{a fortiori} to other
commodities, which add supply elasticity failures (P5) on top of gold's
structural issues. Gold has served as neutral settlement for millennia.
It fails on two properties: P4 on settlement friction and P6 on
seizure. Physical gold settlement incurs 3--8\% costs including
insurance, transport, assay, and storage (BullionVault, 2025). The
efficiency gap versus Bitcoin settlement is approximately four orders
of magnitude at peak congestion and seven orders of magnitude at
prevailing rates, reflecting the fundamental difference between
physical and informational assets. Executive Order 6102 (1933)
demonstrated that gold is seizable at national scale.

\noindent\emph{The verification game.} Gold settlement embeds a verification game
that has no Bitcoin analogue. Model the interaction between a seller
($S$) delivering an asset claimed to be gold and a buyer ($B$)
who must decide whether to verify:

\begin{table}[h]
\centering
\begin{tabular}{lll}
\toprule
& $S$: Authentic & $S$: Forge \\
\midrule
$B$: Trust & $(V - P,\; P)$ & $(-V,\; V + P)$ \\
$B$: Verify & $(V - P - C_V,\; P)$ & $(-C_V,\; -C_F)$ \\
\bottomrule
\end{tabular}
\end{table}

where $V$ is the asset value, $P$ is the price, $C_V$ is
verification cost, and $C_F$ is forgery cost. The equilibrium
depends on the relationship between $C_V$ and $C_F$.

For gold, $C_V$ is \textbf{monotone increasing in forgery
sophistication}: spray-painted lead is caught by visual inspection
(\textasciitilde\$0), but tungsten-core bars require destructive assay
(2-5\% of bar value) or ultrasound testing (\$200-500/bar). Documented
cases include tungsten-core bars in commercial LBMA inventories and
spray-painted lead bars in retail markets. The US gold reserve at Fort
Knox has not been independently audited since 1953, evidence that even
sovereign-scale verification is prohibitively costly.

\begin{proposition}[Verification Cost Asymmetry]
Gold verification cost is strictly increasing in counterfeiter
sophistication, while Bitcoin verification cost is asymptotically zero
and invariant in attacker sophistication.
\end{proposition}

\noindent Let $C_V^\text{gold}(s)$ be the cost of verifying gold against
a counterfeiter of sophistication $s \in [0, 1]$, and let $C_V^\text{btc}$
be the cost of verifying a Bitcoin UTXO. Three claims:

\begin{enumerate}[label=(\roman*),nosep]
\item $(C_V^\text{gold})'(s) > 0$: gold verification cost increases
  with forgery sophistication.
\item $C_V^\text{btc} = \varepsilon \approx 0$, independent of attacker
  sophistication, under Hash (2026a, Assumption 3).
\item Gold verification is non-persistent: each transfer resets the
  verification game, since the bar could have been swapped between
  inspections. Bitcoin verification is persistent: a confirmed UTXO
  remains valid until spent.
\end{enumerate}

\begin{proof}
(i) follows from the physical structure of gold: detecting surface-level
visual fakes costs less than detecting deep fakes via destructive
assay. More sophisticated forgeries require more invasive and more
expensive testing. (ii) follows from Hash (2026a, Assumption 3): verifying a digital
signature
is computationally trivial regardless of the attacker's resources (the
attacker cannot produce a valid signature without the private key).
(iii) follows from the difference between physical and informational
assets: a physical bar's composition can change between inspections; a
UTXO's validity is determined by the blockchain state, which is globally
consistent and tamper-evident.
\end{proof}

The implication for autonomous agents is categorical: an AI agent
verifying gold must either trust a human intermediary, which introduces
counterparty risk at $r = 0$, or deploy physical inspection
infrastructure whose calibration itself depends on human trust chains.
Bitcoin verification requires only a full node, software the agent
controls entirely.

\noindent\emph{The gold-ETF objection.} Gold ETFs settle digitally at low
cost, arguably satisfying P4. But gold ETFs are claims on custodial
gold, not gold itself. They introduce counterparty risk via the
custodian and regulatory capture via the issuer. An ETF can be frozen,
diluted, or restructured by its sponsor. If we accept ETF wrappers as
satisfying P4, then by the same logic Bitcoin ETFs satisfy P4, and the
underlying asset remains the object of analysis, not the wrapper.
Gold's P4 and P6 failures are properties of gold-as-settlement, not
gold-as-financial-product.

\noindent\textbf{C6: Alternative Crypto Assets.} This bucket covers
non-Bitcoin blockchain assets, which divide into two subcategories.

\noindent\textbf{C6a: Alternative L1 tokens.} Every alternative
blockchain protocol has at least one of: a foundation treasury that
constitutes a governance capture surface, a venture capital allocation
that creates concentrated voting power, or a proof-of-stake mechanism
in which wealth concentration maps directly to protocol control.
Ethereum's transition to proof-of-stake in September 2022 made this
explicit: the largest stakers have proportional influence over block
production and transaction ordering.

\noindent\textbf{C6b: Stablecoins.} Tokenized fiat and Treasury claims
settle digitally at low cost and denominate in familiar units, arguably
satisfying P4. But stablecoins fail P2 and P5 by construction. The issuer maintains a freeze function: Tether has frozen
hundreds of addresses at the direction of law enforcement agencies,
totaling over \$1 billion as of 2025. This means
$|CS_\text{protocol,\,stablecoin}| > 0$: a single entity can unilaterally alter
settlement outcomes for any address. Stablecoins also fail P5: they are
denominated in fiat currency, so the underlying asset debases when the
reference currency debases, and the holder bears debasement risk with
none of the yield. The capture is not hypothetical: proposed US
stablecoin legislation (2025--2026) would cap yield payments to holders
below market rates, a provision secured through banking industry
lobbying to protect deposit bases. This is the P2 failure mode in its
clearest form: a political coalition altering the asset's rules to
serve incumbent interests, precisely the capture that neutral
settlement requires immunity from.

\noindent\textbf{C7: Bitcoin.} Bitcoin satisfies all seven properties at
the protocol layer. P1 is supported by the mining Nash equilibrium
(Biais et al., 2019). P2 holds because there is no issuer, no
foundation, and $|CS_\text{protocol}| = 0$. P3 is demonstrated by the
network's continued operation across every jurisdiction that has
attempted to ban it. P4 is empirically verified: \$1B settles for
under \$500 even at peak congestion. P5 is algorithmic: a 21 million
cap enforced by node consensus. P6 follows from cryptographic custody.
P7 is demonstrated by the soft fork upgrade path without governance
capture.

Application-layer capture surfaces exist and should be named directly.
Spot ETF wrappers hold approximately 6\% of transferable supply;
public-company treasuries hold approximately 4\%; exchange reserves
hold approximately 8\%; miner balances are negligible at the supply
scale. Each captured fraction is unilaterally exitable: holders can
withdraw ETF shares into spot Bitcoin and then into self-custody,
miners can switch pools or solo-mine, and node operators can reject
invalid chains. More importantly, under Proposition 3 these fractions
are not load-bearing. Even if ETFs, treasuries, exchanges, and
sovereign reserves in aggregate exceeded 51\% of transferable supply,
protocol neutrality would still hold. Holding coins does not confer
the right to alter the supply schedule, censor transactions, seize
other holders' keys, or change consensus rules. The 21-million cap is
enforced by node operators running the reference implementation, not
by coin holders voting their balances. This is the structural
distinction between Bitcoin and every asset it competes with as
reserve settlement.

\subsection{3.4 Coexistence}

The elimination does not require fiat collapse. Bitcoin operates as
superior collateral and settlement finality in parallel with fiat
payment rails. The cascade pressure (Hash, 2026a, Proposition 2) applies
to the reserve settlement function. Fiat retains payment and
unit-of-account roles, potentially indefinitely. This is not a
revolutionary claim; it is a structural one.
\section{4. Empirical Support}

Theory without data is speculation. The following observations are
consistent with the framework.

\subsection{4.1 Institutional Adoption}

Spot Bitcoin ETFs accumulated \$61.5 billion in net inflows since
January 2024, reaching approximately \$170 billion in total AUM.
BlackRock's IBIT crossed \$97 billion in 435 days, faster than any ETF
in history (The Block, 2025). This is consistent with the cascade
dynamics predicted by the Exit Game: institutional actors crossing their
adoption thresholds $p_i^*$ in clusters.

\subsection{4.2 Sovereign Behavior}

El Salvador holds 7,529 BTC with daily purchases and geothermal mining
(Bitcoin Treasuries, 2025). Singapore issued 13 digital payment token
licenses in 2024 (Monetary Authority of Singapore, 2024). Argentina
received \$91.1 billion in cryptocurrency value amid 143\% inflation
(Chainalysis, 2024). These observations are consistent with sovereign
defection from the Stay coalition (Hash, 2026a, Theorem 2).

\subsection{4.3 Settlement Efficiency}

Bitcoin settles \$1 billion in value for fees under \$1 at prevailing
mempool rates (1 sat/vB, April 2026) and under \$520 at the highest
sustained congestion on record (Mempool.space, 2024--2026), with
mathematical finality in under 60 minutes. Gold settlement incurs costs
across two dimensions that Bitcoin avoids.

\noindent\emph{One-time physical transport} of LBMA-grade bullion at
wholesale costs 0.5--1\% of value for insured armored delivery
(BullionVault, 2025), approximately \$5M--\$10M on \$1B, before assay
at \$200--500 per bar across roughly 1,240 bars at 400 oz each, and
before vault setup. Retail and cross-border settlement incurs 3--8\%.

\noindent\emph{Ongoing ETF custody}, the only format in which gold
settles digitally, runs 17--40 basis points annually (SGOL 0.17\%, IAU
0.25\%, GLD 0.40\%), or roughly \$2.5M per year on \$1B at the
institutional median. Bitcoin self-custody has no recurring protocol
cost.

At wholesale peak congestion, the one-time settlement ratio is
approximately four orders of magnitude (\$5M / \$520 $\approx$
10,000$\times$); at prevailing rates it exceeds seven orders of
magnitude. The recurring custody comparison has no Bitcoin analogue:
ETF custody is a permanent fee stream that also introduces counterparty
surface (P2 failure) and jurisdictional seizure risk (P6 failure). This
is the empirical basis for the P4 elimination of gold.

\subsection{4.4 Network Resilience}

China's comprehensive mining ban in 2021 resulted in total hashrate
recovery to all-time highs within six months (Cambridge Centre for
Alternative Finance, 2022). Bitcoin's hashrate reached approximately
700 EH/s in 2025, with no PoW competitor exceeding 1\% of this security
budget (Cambridge Centre for Alternative Finance, 2025).

\subsection{4.5 Historical Uniqueness}

Bitcoin's position derives from an unreplicable historical sequence.
First, pure origin: no ICO, no pre-mine, no venture capital, no
foundation treasury (Nakamoto, 2008). Second, founderless development:
the creator departed, with an estimated 1.1 million BTC unspent since
2009 (Lerner, 2019). Third, survival of six categories of existential
crisis without protocol failure. Fourth, PoW hashrate dominance
following Ethereum's migration to proof-of-stake in 2022.

No competitor can replicate this sequence. As Huberman, Leshno, and
Moallemi (2021) observe, Bitcoin is a ``monopoly without a monopolist,''
and the monopoly derives from the \emph{absence} of a monopolist, not
despite it. A new protocol with a known founder, a foundation treasury,
and venture backing fails P2 by construction, regardless of its
technical merits.
\section{5. Discussion}

\subsection{5.1 The Gold Question}

Gold is the strongest competitor and deserves extended treatment. Gold
has served as neutral settlement for millennia. Its failure on P4 and P6
is not a deficiency of gold but a consequence of its physicality. The
information revolution created a new possibility: settlement that is
mathematical rather than physical, where custody requires knowledge
rather than location.

Krugman (2018) argues that Bitcoin lacks fundamental value because it
has no tether to economic activity. This is reasonable if you weight
intrinsic yield over settlement properties. The present analysis
weights structure, and on structure, gold fails two properties that
Bitcoin satisfies.

Taleb (2021) makes a different case: that Bitcoin is fragile because it
depends on continuous mining, electricity, and internet infrastructure.
This is the steelman for P7 failure. Our response: seventeen years of
operation through six categories of existential crisis (exchange
failures, regulatory bans, mining exodus, protocol disputes, market
crashes, and sustained negative press) constitute empirical evidence for
P7, though not proof. The fee-transition risk (Carlsten et al., 2016)
remains the strongest version of Taleb's concern.

\subsection{5.2 Attack Survival}

The elimination holds only if Bitcoin actually satisfies P1--P7. The most
serious challenge is sovereign attacks: can state actors ban, seize, or
suppress Bitcoin?

The steelman: a coordinated coalition of G7 plus China simultaneously
criminalizing possession, enforcing through ISP-level filtering, and
imposing secondary sanctions. The defense rests on three structural
features. Informational security (P6) transforms seizure into a
key-management problem. Permissionless access (P3) redirects activity
rather than eliminating it, as China's 2021 ban empirically
demonstrated. And the coalition itself is unstable: the first defector
captures fleeing capital, and coordination problems are precisely what
this paper series is about.

We regard F1, a permanent single global coordinator, as the most
plausible falsification path. Not because it is likely, but because it
is the only attack that does not face a game-theoretic self-undermining
dynamic.

\subsection{5.3 The Volatility Objection}

Bitcoin's annualized volatility has exceeded 70\% in multiple calendar
years, with drawdowns exceeding 50\% from all-time highs occurring in
2014, 2018, 2022, and again in early 2026 at approximately 52\% decline
from cycle highs. This is the most common objection from traditional
finance practitioners and deserves direct treatment.

Volatility is not among P1--P7. This is deliberate: volatility is a
market phenomenon, not a structural property of the protocol. The Exit
Game model (Hash, 2026a, condition M2) posits that
$\sigma_B'(p) < 0$: deeper markets reduce volatility. The long-term
trend is consistent with this: Bitcoin's 90-day realized volatility
has declined from roughly 150\% in 2011 to roughly 45\% in 2024.
However, the empirical evidence is mixed. Sharp drawdowns persist at
current adoption levels. The power law channel, under which price
follows a power law of time since genesis, describes the central
tendency but does not constrain short-term deviations.

We therefore distinguish two claims. The \emph{structural} claim:
volatility affects adoption timing through the $\lambda_i \cdot
\sigma_B(p)$ term but does not affect equilibrium structure. More
risk-averse actors have higher thresholds $p_i^*$ and wait longer. This
holds regardless of $\sigma_B$'s sign. The \emph{empirical} claim:
volatility will decrease with adoption as markets deepen. This is
plausible but unproven at current adoption levels, and the 2026
drawdown is a direct counterexample to the naive version. M2 should be
understood as a long-run tendency, not a monotone empirical fact at
every time horizon.

\subsection{5.4 Settlement vs.~Acceptance}

The elimination framework rests on a distinction that deserves explicit
treatment: \emph{settlement} is not \emph{acceptance}.

Settlement in Bitcoin is a cryptographic fact. A valid signature
transfers value from address A to address B. The protocol does not ask
who owns the addresses, whether the transaction serves a lawful purpose,
or whether the receiving party will be able to spend the output. At the
protocol layer, P3 (permissionless access) holds unconditionally.

Acceptance is a social fact. If address B is associated with a
sanctioned entity, a designation maintained by OFAC, the EU, or other
authorities, then exchanges and custodians will refuse to credit
incoming funds from B or any address that received value from B. The
coins are settled but economically impaired: their acceptance by
regulated counterparties depends on off-chain arbitration that the
protocol cannot control.

This creates a race condition. Mixing services, coinjoins, and
cross-chain bridges attempt to break the provenance chain. Compliance
tools (Chainalysis, Elliptic) attempt to trace it. The effectiveness of
both evolves continuously. Partial laundering creates a spectrum of
``cleanliness'' rather than a binary clean/tainted classification, and
gray-market exchanges accept coins that regulated exchanges reject.

For our analysis, the implication is: P3 guarantees settlement at the
protocol layer, but the \emph{economic utility} of that settlement
depends on the post-settlement acceptance environment. The elimination
holds for settlement functionality, since Bitcoin is the unique asset
where the \emph{protocol} imposes no barriers. Whether the
\emph{ecosystem} around the protocol adds barriers is a compliance
question, not a settlement one. This distinction is sharpened in Hash
(2026c), where autonomous agents face the compliance problem in its
purest form.

\noindent\textbf{The Acceptance Game.} The settlement-acceptance distinction
produces a natural enforcement equilibrium without requiring
protocol-level censorship. Consider an actor receiving coins of
uncertain provenance. The expected payoff of blind acceptance is:

$\E[\text{Accept}] = V - p \cdot d(t) \cdot c$

where $V$ is the value of the coins, $p$ is the probability of
taint, $d(t)$ is the probability of detection at time $t$, and $c$
is the cost of punishment (legal penalties, reputational damage,
asset seizure). Acceptance is rational when $V > p \cdot d(t) \cdot c$.

Three features distinguish this game from the cash analogue. First,
$d(t)$ is monotonically increasing: Bitcoin's permanent
ledger means chain analysis capabilities improve over time and apply
retroactively to past transactions. Coins accepted today may be flagged
in 2030. Cash has no such property: once spent, provenance is lost.
Second, regulated entities such as exchanges and custodians face $c$
values that dominate $V$ for any realistic taint probability, with
license revocation, criminal liability, and institutional collapse on
the downside, and these entities control the fiat on/off ramps. Third,
the resulting ``provenance discount'' on tainted coins reduces the
economic payoff of the underlying crime: you can steal Bitcoin, but
you cannot easily convert it at par value.

The equilibrium stratifies by actor type:

\begin{table}[h]
\centering
\begin{tabular}{llll}
\toprule
Actor & $d$ & $c$ & Filtering \\
\midrule
Regulated exchange & High & Very high & Mandatory \\
Individual (KYC'd) & Medium & High & Incentivized \\
AI agent & Medium & Low & Rational \\
Gray market & Low & Low & Minimal \\
\bottomrule
\end{tabular}
\end{table}

Notably, the actors most capable of filtering (regulated exchanges) are
also the actors most severely punished for failure, creating robust
enforcement at the fiat conversion bottleneck. For autonomous agents,
filtering converges to near-perfection not from moral obligation but
from the trivial cost of on-chain provenance analysis relative to any
nonzero expected punishment.

The policy implication: protocol-level censorship, which would violate
P2 and P3, is unnecessary for enforcement, because the application
layer produces its own enforcement equilibrium through rational
acceptance decisions. Bitcoin's settlement-layer neutrality is
compatible with discretionary filtering at the acceptance layer, and
indeed is strengthened by it.

\subsection{5.5 The Eighth-Category Objection}

The elimination tests seven asset classes. A reader may ask: what about
art, physical commodities, intellectual property, collectibles, carbon
credits, private equity, tokenized real-world assets, or a future
AGI-launched settlement protocol? The answer is reduction: each
candidate collapses structurally into C1--C6.

Art, collectibles, wine, and watches are C4-family: physical, seizable,
assay-dependent, illiquid. Physical commodities such as oil, wheat, and
industrial metals are C5-family with one structural worsening: supply
is elastic ($dS/dP > 0$) because higher prices draw out more
production, so they additionally fail P5. Patents, trademarks,
copyrights, carbon credits, and regulatory allowances are C1-family:
the state is both the issuer and the enforcer, making
$|CS_\text{protocol}| > 0$ by construction, with no exit path since
protocol authority and ownership are fused. Private equity and venture stakes are C3-family.
Crypto-collateralized stablecoins and governance-token protocols are
C6: token-weighted votes over collateral, parameters, or issuance are
foundation-equivalents. Tokenized real-world assets inherit the
underlying's failures plus an issuer-wrapper P2 failure.

The sharpest form of the objection is an AGI-launched successor
protocol. The game-theoretic answer is that a rational coordinator,
whether human or machine, selects the protocol that already satisfies
P1--P7. Proof-of-work Lindy, absence of a known founder, fair emission
history, small block size preserving cheap full-node verification, and
seventeen years of adversarial survival are path-dependent credentials
that a freshly launched protocol cannot replicate. A new protocol with
identifiable founders or a pre-mine fails P2 at genesis; one without
them inherits Bitcoin's failure modes with none of its history. AGI
actors face the same payoff structure as human actors in Hash (2026a):
use what works, and fund its security through otherwise-stranded
energy. Launching a competing protocol reintroduces the coordination
failure the protocol was adopted to escape. Hash (2026c) develops the
full game theory of Bitcoin neutral settlement under autonomous-agent
adoption.

Even granting an unreduced edge case, the reserve-settlement function
requires stock sufficient to absorb sovereign-scale flows without
price dislocation that would destroy the reserve property itself. Art,
collectibles, and intellectual-property markets are orders of magnitude
smaller in transferable stock than the aggregate reserve base they
would need to substitute for. The elimination is therefore robust to
the eighth-category objection: the absent category either reduces to
an already-failed bucket, or is irrelevant at the scale the question
concerns.

\subsection{5.6 Limitations}

The elimination is as strong as the property framework. If P1--P7 are
incomplete, meaning there exists an eighth necessary property that
Bitcoin violates, the elimination fails. The sufficiency argument
(Proposition 2) addresses this, but sufficiency proofs are inherently
harder than necessity proofs because they require exhaustive
enumeration of alternatives.

\noindent\emph{On minimality and completeness.} We claim that P1--P7
are individually necessary and jointly sufficient. We do not claim
they are \emph{minimal}, meaning that no proper subset of P1--P7
suffices. Formal minimality would require showing that for each
property $P_k$, there exists an asset satisfying all $P_j$ for $j \neq
k$ that fails as neutral settlement. The counterexamples are suggestive
(Ethereum for P1, Solana for P2, CBDCs for P3, gold for P4, commodity
bonds for P5, Swiss banking for P6, Litecoin for P7), but a rigorous
minimality proof in the tradition of axiomatic characterizations in
social choice (Arrow, 1951) and mechanism design (Gibbard, 1973) would
require 7--10 additional pages and is not attempted here. The relevant
question for the elimination is necessity plus sufficiency: can an
asset fail any single property and still serve as neutral settlement?
The answer is no, and that is what the framework proves.

Seventeen years of data is brief for a monetary claim. Gold has
millennia. The response, that the relevant properties are structural
and not historical, is logically sound but empirically untested at the
timescales that matter. Carlsten et al.'s (2016) analysis of the
fee-only security budget is particularly relevant: the transition from
block subsidies to fee-only security is the largest untested structural
assumption, and it will not be resolved by theory alone.
\section{6. Conclusion}

Global wealth is held in seven asset buckets: fiat currencies,
sovereign bonds, equities, real estate, gold, alternative crypto assets
(alt-L1 tokens and stablecoins), and Bitcoin. Neutral settlement
requires seven properties: protocol security, neutrality, permissionless
access, cheap finality, absolute scarcity, informational security, and
adaptive resilience. The two sevens are independent. One partitions
where capital sits, the other specifies what reserve settlement
requires.

Six of the seven buckets fail at least one of the seven properties,
and the failures are constitutive rather than contingent: a fiat
currency without a discretionary issuer is not a fiat currency; gold
without physicality is not gold; a smart-contract platform without a
founder or treasury did not historically happen. Bitcoin fails none at
the protocol layer. At the application layer, it is the unique asset
whose captured fractions cannot translate into protocol control.

The claim is not that Bitcoin dominates on any single dimension. Gold
is older, equities yield more, fiat is more liquid for daily payments.
Nor is the claim that Bitcoin is uncapturable. ETF wrappers, exchange
custody, corporate treasuries, and mining pool concentration are real
capture surfaces, and each continues to grow. Ownership concentration
may one day exceed 51\% across institutional allocators. The claim is that
none of this alters protocol neutrality, because under proof-of-work
with cryptographic custody, holding coins confers no consensus vote
and no ability to prevent other key-holders from settling. This is the
structural difference Proposition 3 formalizes, and it is the
difference that Hash (2026a) takes as given without proving. Hash
(2026a) establishes why capital flows toward neutral settlement; this
paper identifies which asset satisfies neutrality under partial
capture; Hash (2026c) and Hash (2026d) treat the autonomous-agent and
multipolar-sovereign regimes in which that flow accelerates.
\section{References}

\begin{description}[style=nextline,leftmargin=0pt,labelsep=0pt,font=\normalfont]

\item Biais, B., Bisiere, C., Bouvard, M., and Casamatta, C.\ (2019). The
blockchain folk theorem. \emph{Review of Financial Studies}, 32(5), 1662--1715.
\url{https://doi.org/10.1093/rfs/hhy095}

\item Bitcoin Treasuries (2025). Sovereign Bitcoin holdings tracker.
bitcointreasuries.net.

\item Budish, E.\ (2018). The economic limits of Bitcoin and the blockchain.
\emph{NBER Working Paper} No.~24717.
\url{https://doi.org/10.3386/w24717}

\item BullionVault (2025). Gold settlement cost analysis.

\item Cambridge Centre for Alternative Finance (2022). Cambridge Bitcoin
Electricity Consumption Index: China mining ban impact.

\item Cambridge Centre for Alternative Finance (2025). Cambridge Bitcoin
Electricity Consumption Index: Global hashrate data.

\item Carlsten, M., Kalodner, H., Weinberg, S.\ M., and Narayanan, A.\ (2016).
On the instability of Bitcoin without the block reward.
\emph{Proceedings of the 2016 ACM SIGSAC}, 154--167.
\url{https://doi.org/10.1145/2976749.2978408}

\item Chainalysis (2024). The 2024 Geography of Cryptocurrency Report.

\item Federal Reserve Bank of St.~Louis (2024). FRED Economic Data: 10-Year
Treasury minus inflation expectations (DGS10, T10YIE series).

\item Hash (2026a). Bitcoin exit dominance in monetary coordination games.
SSRN Working Paper 6299081. \url{https://ssrn.com/abstract=6299081}.

\item Hash (2026c). Settlement at zero trust: Bitcoin and autonomous economic
agents. Working paper, \url{bitcoingametheory.com}.

\item Hash (2026d). Monetary predator-prey dynamics: Enforcement gridlock
and neutral settlement survival. Working paper, \url{bitcoingametheory.com}.

\item Huberman, G., Leshno, J.\ D., and Moallemi, C.\ (2021). Monopoly without a
monopolist. \emph{Review of Economic Studies}, 88(6), 3011--3040.
\url{https://doi.org/10.1093/restud/rdab014}

\item Krugman, P.\ (2018). Transaction costs and tethers: Why I'm a crypto
skeptic. \emph{The New York Times}, July 31.

\item Lerner, S.\ (2019). The return of the deniers and the revenge of Patoshi.

\item Mempool.space (2024). Bitcoin transaction fee and settlement data.

\item Monetary Authority of Singapore (2024). Digital Payment Token Licensing
Statistics.

\item Nakamoto, S.\ (2008). Bitcoin: A peer-to-peer electronic cash system.

\item Taleb, N.\ N.\ (2021). Bitcoin, currencies, and fragility.
\emph{Quantitative Finance}, 21(8), 1249--1255.
\url{https://doi.org/10.1080/14697688.2021.1952702}

\item The Block (2025). Bitcoin ETF AUM tracker.

\end{description}

% ============================================================
% APPENDIX
% ============================================================

\appendix

\section{Notation}

\begin{table}[h]
\centering
\begin{tabularx}{\textwidth}{lX}
\toprule
Symbol & Definition \\
\midrule
P1--P7 & Seven necessary properties: Protocol Security, Neutrality, Permissionless Access, Cheap Finality, Absolute Scarcity, Informational Security, Adaptive Resilience \\
$|CS_\text{protocol}|$ & Protocol-layer capture surface: parameters an entity can unilaterally alter without network consensus \\
$\kappa_\text{protocol}(A)$ & Fraction of $A$'s settlement rules, supply schedule, or transaction validity any feasible coalition can alter \\
$\kappa_\text{owner}(A)$ & Maximum fraction of $A$'s transferable supply any coalition can hold, seize, freeze, or immobilize \\
$dS/dP$ & Supply elasticity ($dS/dP = 0$ under absolute scarcity) \\
$C_\text{attack}$ & Cost of attacking the protocol \\
$P_\text{collapse}$ & Probability of network collapse conditional on attack \\
$\Delta_\text{attack}$ & Attacker's gain from a successful attack \\
$\Pi_\text{honest}$ & Per-period profit from honest mining \\
$\delta$ & Discount factor \\
$V, P$ & Asset value, transaction price (Verification Game, \S3.3 C5) \\
$C_V, C_F$ & Verification cost, forgery cost (Verification Game) \\
$s \in [0, 1]$ & Counterfeiter sophistication \\
$\varepsilon$ & Near-zero computational verification cost for Bitcoin UTXOs \\
$d(t)$ & Probability of taint detection at time $t$ (Acceptance Game, \S5.4) \\
$p, c$ & Taint probability, punishment cost (Acceptance Game) \\
$\lambda_i$ & Actor $i$'s risk aversion \\
$\sigma_B(p)$ & Volatility of neutral settlement asset at adoption level $p$ \\
$\sigma_B'(p) < 0$ & Volatility decreases as adoption deepens (Condition M2) \\
$p_i^*$ & Actor $i$'s adoption threshold \\
\bottomrule
\end{tabularx}
\end{table}

\end{document}
